\subsection{典范自由分解}

对每个 A-模 M，用 \(\mathbf{L}_{0}(\mathbf{M})\) 表示以 M 为基的自由 A-模 \(\mathbf{A}^{(\mathbf{M})}\) ， \((e_{m})_{m\in\mathbf{M}}\) 为其典范基， \(p_{\mathbf{M}}:\mathbf{L}_{0}(\mathbf{M})\to\mathbf{M}\) 为使得

\[
p _ {\mathbf {M}} (e _ {m}) = m, \quad m \in \mathbf {M}.
\]

令 \(Z_0(\mathbf{M}) = \operatorname{Ker} p_{\mathbf{M}}\) 并设 \(i_{\mathbf{M}}: Z_0(\mathbf{M}) \to L_0(\mathbf{M})\) 为典范嵌入。我们有一个正合列

\[
0 \longrightarrow Z _ {0} (\mathrm{M}) \stackrel {i _ {\mathrm{M}}} {\longrightarrow} L _ {0} (\mathrm{M}) \stackrel {p _ {\mathrm{M}}} {\longrightarrow} \mathrm{M} \longrightarrow 0.\tag{1}
\]

我们通过如下设定来定义分次模 L(M)： \(\mathrm{L}_{n}(\mathrm{M}) = 0\) （当 n \textless{} 0），并对整数 n \textgreater{} 0 作归纳

\[
\mathrm{L} _ {n} (\mathrm{M}) = \mathrm{L} _ {0} \left(\mathrm{Z} _ {n - 1} (\mathrm{M})\right); \quad \mathrm{Z} _ {n} (\mathrm{M}) = \mathrm{Z} _ {0} \left(\mathrm{Z} _ {n - 1} (\mathrm{M})\right).\tag{2}
\]

我们定义 A-同态 \(d_{n}^{\mathbf{M}}:\mathrm{L}_{n}(\mathbf{M})\to \mathrm{L}_{n - 1}(\mathbf{M})\) 如下

\[
\left\{ \begin{array}{l l} d _ {n} ^ {\mathbf {M}} = 0, & n \leqslant 0, \\ d _ {1} ^ {\mathbf {M}} = i _ {\mathbf {M}} \circ p _ {\mathbf {Z} _ {0} (\mathbf {M})}, & \\ d _ {n} ^ {\mathbf {M}} = i _ {\mathbf {Z} _ {n - 2} (\mathbf {M})} \circ p _ {\mathbf {Z} _ {n - 1} (\mathbf {M})}, & n > 1. \end{array} \right.\tag{3}
\]

由构造，我们有一个正合列

\[
\longrightarrow \mathrm{L} _ {n} (\mathrm{M}) \xrightarrow {d _ {n} ^ {\mathrm{M}}} \mathrm{L} _ {n - 1} (\mathrm{M}) \longrightarrow \dots \longrightarrow \mathrm{L} _ {0} (\mathrm{M}) \xrightarrow {p _ {\mathrm{M}}} \mathrm{M} \longrightarrow 0,
\]

因此，如果我们将 \(p_{\mathbf{M}}\) 作为复形态射进行延拓

\[
p _ {\mathbf {M}}: (\mathrm{L} (\mathbf {M}), d ^ {\mathbf {M}}) \rightarrow \mathbf {M},
\]

得到M的一个自由分解，称为M的典范自由分解。

设 \(f: \mathbf{M} \to \mathbf{N}\) 是 A-模的同态。记

\[
\mathrm{L} _ {0} (f): \mathrm{L} _ {0} (\mathrm{M}) \rightarrow \mathrm{L} _ {0} (\mathrm{N})
\]

为唯一的 A-同态，使得 \(L_0(f)(e_m) = e_{f(m)}\) 对所有 \(m \in \mathbf{M}\) 。我们有

\[
p _ {\mathbf {N}} \circ \mathrm{L} _ {0} (f) = f \circ p _ {\mathbf {M}}.\tag{4}
\]

因此， \(L_0(f)\) 诱导一个A-同态 \(Z_0(f):Z_0(M)\to Z_0(N)\) 并且我们有

\[
i _ {\mathbf {N}} \circ \mathbf {Z} _ {0} (f) = \mathbf {L} _ {0} (f) \circ i _ {\mathbf {M}}.\tag{5}
\]

令 \(\mathrm{L}_{n}(f)=0\) （对 n\textless0），并在整数 n\textgreater0 上归纳地定义同态 \(\mathrm{L}_{n}(f):\mathrm{L}_{n}(\mathrm{M})\to\mathrm{L}_{n}(\mathrm{N})\) 与 \(\mathrm{Z}_{n}(f):\mathrm{Z}_{n}(\mathrm{M})\to\mathrm{Z}_{n}(\mathrm{N})\) 如下

\[
\left\{ \begin{array}{l} \mathrm{L} _ {n} (f) = \mathrm{L} _ {0} (\mathrm{Z} _ {n - 1} (f)) \\ \mathrm{Z} _ {n} (f) = \mathrm{Z} _ {0} (\mathrm{Z} _ {n - 1} (f)). \end{array} \right.\tag{6}
\]

命题 4. --- L(f) : L(M) → L(N) 是 A-模复形的一个态射；此外，我们有 \(p_{\mathbf{M}} \circ \mathrm{L}(f) = f \circ p_{\mathbf{N}}\) .

我们需要证明，对于每个整数 n \textgreater{} 0，公式

\[
d _ {n} ^ {\mathbf {N}} \circ \mathrm{L} _ {n} (f) = \mathrm{L} _ {n - 1} (f) \circ d _ {n} ^ {\mathbf {M}}.
\]

我们首先有

\[
\begin{array}{r l r l} d _ {1} ^ {\mathbf {N}} \circ \mathrm{L} _ {1} (f) & = i _ {\mathrm{N}} \circ p _ {\mathbf {Z} _ {0} (\mathbf {N})} \circ \mathrm{L} _ {0} (\mathbf {Z} _ {0} (f)) & & (\mathrm {by~(3)~and~(6)}) \\ & = i _ {\mathrm{N}} \circ \mathrm{Z} _ {0} (f) \circ p _ {\mathbf {Z} _ {0} (\mathbf {M})} & & (\mathrm {by~(4)}) \\ & = \mathrm{L} _ {0} (f) \circ i _ {\mathbf {M}} \circ p _ {\mathbf {Z} _ {0} (\mathbf {M})} & & (\mathrm {by~(5)}) \\ & = \mathrm{L} _ {0} (f) \circ d _ {1} ^ {\mathbf {M}} & & (\mathrm {by~(3))}. \end{array}
\]

当 \(n > 1\) 时，我们依次有

\[
\begin{array}{l l} d _ {n} ^ {\mathbf {N}} \circ \mathrm{L} _ {n} (f) = i _ {\mathrm{Z} _ {n - 2 (\mathrm{N})}} \circ p _ {\mathrm{Z} _ {n - 1 (\mathrm{N})}} \circ \mathrm{L} _ {0} (\mathrm{Z} _ {n - 1} (f)) & \text {(by (3) and (6))} \\ = i _ {\mathrm{Z} _ {n - 2 (\mathrm{N})}} \circ \mathrm{Z} _ {n - 1} (f) \circ p _ {\mathrm{Z} _ {n - 1} (\mathrm{M})} & \text {(by (4))} \\ = i _ {\mathrm{Z} _ {n - 2 (\mathrm{N})}} \circ \mathrm{Z} _ {0} (\mathrm{Z} _ {n - 2} (f)) \circ p _ {\mathrm{Z} _ {n - 1} (\mathrm{M})} & \text {(by (6))} \\ = \mathrm{L} _ {0} (\mathrm{Z} _ {n - 2} (f)) \circ i _ {\mathrm{Z} _ {n - 2 (\mathrm{M})}} \circ p _ {\mathrm{Z} _ {n - 1} (\mathrm{M})} & \text {(by (5))} \\ = \mathrm{L} _ {n - 1} (f) \circ d _ {n} ^ {\mathbf {M}} & \text {(by (3) and (6)) .} \end{array}
\]

我们立即有

\[
\mathrm{L} (1 _ {\mathbf {M}}) = 1 _ {\mathrm{L} (\mathbf {M})}.\tag{7}
\]

另一方面，如果 \(g: \mathbf{N} \to \mathbf{P}\) 是A-模的同态，我们有

\[
\mathrm{L} (g \circ f) = \mathrm{L} (g) \circ \mathrm{L} (f).\tag{8}
\]

确实，对于 \(m \in \mathbf{M}\) ,

\[
\mathrm{L} _ {0} (g \circ f) \left(e _ {m}\right) = e _ {g \circ f (m)} = \mathrm{L} _ {0} (g) \left(e _ {f (m)}\right) = \mathrm{L} _ {0} (g) \circ \mathrm{L} _ {0} (f) \left(e _ {m}\right),
\]

因此 \(\mathrm{L}_{0}(g\circ f)=\mathrm{L}(g)\circ\mathrm{L}(f)\) ；因此 \(\mathrm{Z}_{0}(g\circ f)=\mathrm{Z}_{0}(g)\circ\mathrm{Z}_{0}(f)\) ；由此立即得出 \(\mathrm{L}_{n}(g\circ f)=\mathrm{L}_{n}(g)\circ\mathrm{L}_{n}(f)\) ，对于 \(n\geqslant0\) ，通过对 n 进行归纳，由此得到 (8)。

注记。--- 若 f, g ∈ Hom\_A (M, N)，则不一定有 L(f + g) = L(f) + L(g)。然而，根据 X，第 49 页，命题 3，这两个态射是同伦的。

设 M 为一个右 A-模；用 A° 表示 A 的反环，用 M° 表示 M 的底层 A°-模，并用 L(M°) 表示其典范自由分解。我们用 L(M) 表示 L(M°) 的底层 A-复形 L(M°)°，并称其为 M 的典范自由分解。于是我们有

\[
\mathrm{L} (\mathbf {M} ^ {\circ}) = \mathrm{L} (\mathbf {M}) ^ {\circ}.
\]
% semantic-fragments: legacy-input-seam
\relax{}
